% =====================================================================
%  Report 4
%  How the Java code implements the O'FAIRe questions
%  Build: latexmk -pdf emmanuel_dataTerra_Report4-OFAIRe-Java-Implementation.tex
%  Everything here comes from reading the source code of O'FAIRe
%  (github.com/agroportal/fairness, branch master, commit 24def6b of 4 June 2026).
% =====================================================================
\documentclass[11pt,a4paper]{article}
\usepackage{dataterra-report}

\reportsubtitle{How the Java code implements the O'FAIRe questions}

% \band{F1}: the grey line that opens the questions of one sub-principle
\newcommand{\band}[1]{\multicolumn{4}{@{}l}{\cellcolor{dtband}\textbf{\color{dtnavy}#1}}\\*}

\begin{document}
\makereporttitlepage

% =====================================================================
\section{Introduction}

In this short report, our goal is to look at how the Java code of O'FAIRe implements the 61
questions presented in Report~2. The JSON file says what is asked and how many points an answer
is worth, but it is the code that decides which answer an artefact receives. Since the scores
shown in EarthPortal come from this code, I wanted to read it once, question by question.

% =====================================================================
\section{Development}

\subsection{Where the code is, and what I read}

The source code of O'FAIRe is public\footnote{\url{https://github.com/agroportal/fairness}, folder
\prop{src/main/java/fr/lirmm/fairness/assessment}.}. I read it as it was on 4 June 2026 (branch
\prop{master}, commit \prop{24def6b}). 

Each sub-principle has one Java class (\prop{F1.java}, \prop{F2.java}, and so on until
\prop{R13.java}), in the folder \prop{principles/criterion/impl}. For one artefact, the code asks
the portal API for the latest submission with all its fields, and then goes through the questions
one by one. For most questions, the code only checks that the field has a value; I then say
that the field is \emph{filled}. Whether the value is correct is not checked, except where
Table~\ref{tab:rules} says so.

The points of the questions are added for each sub-principle, then for each principle, then for the
artefact. The score out of 100 is the total divided by 478, and the decimals are dropped (for
example, 77.9 is shown as 77).

\subsection{Four kinds of questions}

The code does not treat the 61 questions in the same way. Table~\ref{tab:kinds} gives the four cases.

\begin{center}
\small
\renewcommand{\arraystretch}{1.2}
\begin{tabular}{@{}l>{\raggedright\arraybackslash}p{9.2cm}rr@{}}
\toprule
\textbf{Kind} & \textbf{What the code does} & \textbf{Questions} & \textbf{Points} \\
\midrule
Tested & It reads one or several fields of the artefact and chooses the answer. & 38 & 300 \\
Given & It gives the best answer to every artefact, because the answer depends on the portal
  and not on the artefact. & 12 & 106 \\
Not implemented & It gives 0 to every artefact, with the message \textit{This question cannot be
  evaluated in the repository at this time}. & 10 & 70 \\
Always 0 & F3Q1 only: it gives 0 to every artefact, with an unfinished message. & 1 & 21 \\
\midrule
\textbf{Total} & & \textbf{61} & \textbf{497} \\
\bottomrule
\end{tabular}
\captionof{table}{The four kinds of questions in the code. \textit{Points} is the sum of the best
answers.}
\label{tab:kinds}
\end{center}

\subsection{Question by question}

Table~\ref{tab:rules} gives the rule of each question in one or two lines. The fields are written
with the names the portal uses in its API; Report~3 gives the MOD name of each of them. A
\emph{valid URL} is a text that has the form of an \prop{http} or \prop{https} address. A URL
\emph{resolves} when the connection is made within one second and the server answers with the
code 200, or with the code 302 and a content type; redirections are not followed.

{\footnotesize
\setlength{\tabcolsep}{4pt}
\renewcommand{\arraystretch}{1.2}
\begin{longtable}{@{}p{1.5cm}r>{\raggedright\arraybackslash}p{2.6cm}>{\raggedright\arraybackslash}p{10.45cm}@{}}
\caption{What the code does for each question. \textit{Max} is the best answer of the question.}
\label{tab:rules}\\
\toprule
\textbf{Question} & \textbf{Max} & \textbf{Kind} & \textbf{What the code does} \\
\midrule
\endfirsthead
\toprule
\textbf{Question} & \textbf{Max} & \textbf{Kind} & \textbf{What the code does} \\
\midrule
\endhead
\bottomrule
\endfoot
\band{Findable}
F1Q1 & 9 & Tested & \prop{URI}: empty 0; filled but not a valid URL 3; a valid URL 9. \\
F1Q2 & 11 & Tested & \prop{identifier}: empty 0; not a valid URL 3; a valid URL 6; a
  \prop{doi.org} address that the DOI service does not know 9; a DOI that the DOI service knows 11.
  With several identifiers, the last one decides. \\
F1Q3 & 12 & Given & \\
F1Q4 & 9 & Tested & \prop{versionIRI}: empty 0; not a valid URL 2; a valid URL that does not
  resolve 4; a URL that resolves 9. \\
F2Q1 & 16 & Tested & Counts the filled fields among 25: the 24 of the list in Report~3, and
  \prop{translator}. 2 points for each, up to 8 fields. \\
F2Q2 & 4 & Tested & Counts the filled values among six: the link to the metrics, the numbers of
  classes, properties, individuals and axioms, and \prop{preferredNamespacePrefix}. The points are
  0, 1, 2, 3, 3, 4, 4. \\
F2Q3 & 7 & Tested & Counts the filled fields among the 34 behind the list in Report~3. 1 point for
  each, up to 7 fields. \\
F3Q1 & 21 & Always 0 & The message is \textit{This question}, and nothing more. \\
F3Q2 & 11 & Given & \\
F3Q3 & 10 & Given & \\
F4Q1 & 6 & Tested & Starts at 1 if the portal itself is in the list of repositories of O'FAIRe.
  Then, for each value of \prop{includedInDataCatalog}, adds 1 if it is in the list of repositories
  and 1 if it is in the list of libraries. 2 points for each unit, up to 3. \\
F4Q2 & 10 & Tested & The same count, with the list of repositories only. 2 points for each unit,
  up to 5. \\
F4Q3 & 8 & Not implemented & \\
\band{Accessible}
A1Q1 & 6 & Tested & Tests whether \prop{URI} and each \prop{identifier} is a valid URL that
  resolves. One of them is enough for the 6 points. \\
A1Q2 & 7 & Given & \\
A1Q3 & 24 & Tested & Asks the \prop{URI}, and then the address of the metadata in the portal API,
  for five formats (JSON, XML, HTML, plain text, RDF/XML), and counts the formats accepted. 3 points
  for each, up to 8. \\
A1Q4 & 6 & Tested & \prop{endpoint} filled: 6. The endpoint itself is not tested. \\
A1.1Q1 & 20 & Given & \\
A1.1Q2 & 4 & Given & \\
A1.1Q3 & 4 & Given & \\
A1.2Q1 & 11 & Given & \\
A1.2Q2 & 11 & Given & \\
A2Q1 & 7 & Given & \\
A2Q2 & 5 & Given & \\
A2Q3 & 4 & Given & \\
A2Q4 & 4 & Tested & Counts \prop{status} and \prop{deprecated}. 2 points for each. \\
\band{Interoperable}
I1Q1 & 18 & Tested & \prop{hasOntologyLanguage} must be exactly PDF, TXT, CSV, XML, OBO, RDFS, SKOS
  or OWL, for 5, 7, 9, 10, 11, 12, 15 or 18 points. Any other value: 0. \\
I1Q2 & 10 & Tested & \prop{hasOntologyLanguage} is exactly XML, RDFS, SKOS or OWL: 10. \\
I1Q3 & 5 & Tested & \prop{hasOntologySyntax} filled: 5. \\
I1Q4 & 5 & Tested & \prop{hasFormalityLevel} filled: 5. \\
I1Q5 & 4 & Tested & \prop{hasFormat} and \prop{isFormatOf} both filled: 4. One of them alone: 0. \\
I2Q1 & 5 & Tested & \prop{useImports} filled: 5. \\
I2Q2 & 5 & Tested & \prop{ontologyRelatedTo} filled: 5. \\
I2Q3 & 3 & Not implemented & \\
I2Q4 & 5 & Tested & \prop{isAlignedTo} filled: 5. \\
I2Q5 & 7 & Not implemented & \\
I2Q6 & 2 & Tested & At least one of \prop{similarTo}, \prop{translationOfWork},
  \prop{explanationEvolution}, \prop{generalizes} and \prop{viewOf} filled: 2. \\
I2Q7 & 5 & Tested & O'FAIRe has a list of 34 metadata vocabularies, each with a weight of 0.25, 0.5
  or 0.75. The result is 5 times the average weight of the vocabularies of \prop{metadataVoc} that
  are in this list. \\
I3Q1 & 20 & Not implemented & \\
I3Q2 & 6 & Not implemented & \\
I3Q3 & 7 & Tested & Looks at the fields of the list in Report~3 and at \prop{naturalLanguage}, and
  counts those whose values are all valid URLs. 1 point for each, up to 7. \\
\band{Reusable}
R1Q1 & 5 & Tested & Counts \prop{prefLabelProperty}, \prop{synonymProperty},
  \prop{definitionProperty}, \prop{authorProperty} and \prop{obsoleteProperty}. 1 point for each. \\
R1Q2 & 3 & Tested & Counts \prop{hierarchyProperty}, \prop{obsoleteParent} and the metric
  \prop{maxDepth}. 1 point for each. \\
R1Q3 & 7 & Not implemented & \\
R1Q4 & 6 & Not implemented & \\
R1Q5 & 6 & Not implemented & \\
R1Q6 & 5 & Not implemented & \\
R1.1Q1 & 15 & Tested & \prop{hasLicense}: empty 0; not a valid URL 4; a valid URL that does not
  resolve 8; a URL that resolves but accepts none of the five formats of A1Q3 12; a URL that accepts
  at least one of them 15. \\
R1.1Q2 & 7 & Tested & \prop{viewingRestriction} filled: 4; \prop{morePermissions} filled as well: 7. \\
R1.1Q3 & 15 & Tested & Counts \prop{useGuidelines} and \prop{contact}: 0, 7 or 15. The copyright
  holder of the question is read from the contacts. \\
R1.2Q1 & 8 & Tested & Counts \prop{hasCreator}, \prop{curatedBy}, \prop{hasContributor} and
  \prop{translator}: 0, 3, 5, 7 or 8. \\
R1.2Q2 & 6 & Tested & Counts \prop{source}, \prop{wasGeneratedBy} and \prop{wasInvalidatedBy}.
  2 points for each. \\
R1.2Q3 & 6 & Tested & Counts \prop{accrualMethod}, \prop{accrualPeriodicity} and
  \prop{accrualPolicy}. 2 points for each. \\
R1.2Q4 & 4 & Tested & Counts \prop{version}, \prop{hasPriorVersion} and the link to the list of
  submissions. 2 points for each, up to 2. \\
R1.2Q5 & 2 & Not implemented & \\
R1.2Q6 & 6 & Tested & Counts \prop{usedOntologyEngineeringTool},
  \prop{usedOntologyEngineeringMethodology} and \prop{conformsToKnowledgeRepresentationParadigm}.
  2 points for each. \\
R1.2Q7 & 4 & Tested & Counts \prop{designedForOntologyTask} and \prop{competencyQuestion}.
  2 points for each. \\
R1.2Q8 & 2 & Tested & \prop{fundedBy} filled: 2. \\
R1.3Q1 & 10 & Tested & Counts \prop{endorsedBy} and the projects that use the artefact in the
  portal. 5 points for each. \\
R1.3Q2 & 20 & Tested & Looks at the groups of the artefact: 10 if a group name contains WHEAT, CROP
  or INRAE; 15 if it contains OBO-FOUNDRY; 20 if, in addition, the acronym is BFO, CHEBI, GO, PATO,
  PO or PR. \\
R1.3Q3 & 6 & Tested & \prop{viewingRestriction} contains the word \textit{public}: 6. \\
\end{longtable}}


% =====================================================================
\section{Conclusion}
Throught the reading of the code, I noticed a few points that are not obvious from the JSON file.

\begin{itemize}
  \item \textbf{The lists of properties} in the JSON file does not decide the score, it only decides what is displayed next to the score.
  
  Actually, in \texttt{questions.config.json} (Report 2), most questions carry a properties list. 
  For example, I1Q5 lists \texttt{dct:hasFormat} and \texttt{dct:isFormatOf}.
  Reading this, We would think that the program takes this list and checks those properties to compute the score.
  However, the code works differently,in two separate steps:
 \begin{itemize}
  \item
  \textbf{The test}. 
  Each question has its own small piece of Java, and the fields to check are written by hand inside it.
  For I1Q5, the code says, in substance, "if \texttt{hasFormat} is filled and \texttt{isFormatOf} is filled, give 4 points, otherwise 0". 
  This step never looks at the JSON list.
  \item
  \textbf{The display}. 
  Once the score is chosen, another piece of code (\texttt{Tester.java}) takes the JSON list and
  copies the values of those properties into the result. This is what we see in the O'FAIRe panel, next to the explanation. 
  For \textbf{OSO}, it shows the six formats under \texttt{dct:hasFormat} and \textit{null} under \texttt{dct:isFormatOf}.
  
So there are two lists for each question: one in the Java code, which gives the points, and one in the JSON, which only shows values to the reader.
Most of the time the two lists are the same, but they are maintained separately, and in three questions they have drifted apart:
\begin{itemize}
  \item F2Q1: the code counts the \texttt{translator}, but the JSON list does not show it.
  \item I3Q3: the code also checks the \texttt{naturalLanguage}, which the JSON list does not show. The JSON list shows the subjects, but the code can never count them.
  \item R1.3Q2: the code also looks at the \texttt{acronym}, which is not in the JSON list.
\end{itemize}
\end{itemize}

  \item \textbf{F3.} F3Q1 always gives 0, F3Q2 and F3Q3 are always given, and the weight of F3 (21)
    counts only these two. Every artefact therefore has 21 out of 21 on F3.
  \item \textbf{I1.} The best answers of I1 add up to 42, while its weight is 44. No artefact can
    have more than 95\,\% on I1.
  \item \textbf{I2Q7.} The highest weight in the list of vocabularies is 0.75, so the result cannot
    be more than 3.75 out of 5.
  \item \textbf{A1Q1.} The answer worth 3 points is never used: one identifier that resolves
    already gives 6.
  \item \textbf{R1.3Q2.} The names of the groups are written in the code. None of the groups used
    in EarthPortal has such a name, so its artefacts have 0 out of 20.
  \item \textbf{F4Q1 and F4Q2.} EarthPortal is not in the list of repositories of O'FAIRe. An
    artefact that declares only EarthPortal in \prop{includedInDataCatalog} has 0 on both questions.
\end{itemize}

For an artefact of EarthPortal, this leaves 386 of the 478 points within reach: 478, minus the 70
points of the ten questions that are not implemented, minus the 20 points of R1.3Q2, minus the
2 points that no question of I1 gives.


\end{document}
